% Physical unit 042; accepted Session 03 after fresh 400-dpi authority replay.
\GordanUnitStart{042}{printed page 37}
\GordanTopFolio{37}
\begingroup\small
inest et altero quidem loco tum in $\cos l_0''$ tum in angulis $\varphi_2$ et $\varphi_1$, quos formulis:
\[
 \sin l_0''\cos\varphi_2=\sin l_2,
 \qquad
 \sin l_0''\cos\varphi_1=\sin l_1
\]
definiri notum est.
Attamen, quod
\[
 {e^2\over2}\cos l_0''(\varphi_2-\varphi_1)
\]
minimum est, ea, quae hoc continentur, computari licet sine vitio, loco $l_0''$, jam nota quantitate $l_0'$ posita.

Fieri potest, ut $l_0'$ definitum ex aequatione II valorem, prope $l_0$ situm, non exhibeat, tum quidem, cum summa $l_1+l_2$ minima est et $\psi$ non longe a $\pi$ differt.

Ubi $\tg^2l_0$ repraesentatur, ut facile intelligi potest, quotientem quantitatum minimarum, unde sequitur, ut vel minimo incremento, addendo integrale $\chi$ ad quantitatem $\psi$, adducto, magna valoris $l_0$ mutatio existat.

Quodsi accidit, $l_0$ computandi causa ad ipsam formulam II confugiendum est; quae quia, ponendo loco $l_0''$ valorem prope accedentem $l_0'$ jamque notum, ad elegantiorem formam redigi non potest, alia utendum est ratione scilicet ponendi pro quibusdam quantitatibus alias his proxime sitas.

Quoniam quantitatem
\[
 \psi+{e^2\over2}\cos l_0''(\varphi_2-\varphi_1)
\]
in aequatione III, non multo differentem a $\pi$, possumus, sumere licet:
\[
 \sin\left[\psi+{e^2\over2}\cos l_0''(\varphi_2-\varphi_1)\right]
=\pi-\psi-{e^2\over2}\cos l_0''(\varphi_2-\varphi_1)
\]
\[
 \cos\left[\psi+{e^2\over2}\cos l_0''(\varphi_2-\varphi_1)\right]
=-1+{1\over2}\left\{\pi-\psi-{e^2\over2}\cos l_0''(\varphi_2-\varphi_1)\right\}^2
\]
unde fit:
\[
\begin{aligned}
\tg^2l_0''\left\{\pi-\psi-{e^2\over2}\cos l_0''(\varphi_2-\varphi_1)\right\}^2
&=\tg^2l_1+\tg^2l_2+2\tg l_1\tg l_2\\
&\quad-\tg l_1\tg l_2\left\{\pi-\psi-{e^2\over2}\cos l_0''(\varphi_2-\varphi_1)\right\}^2
\end{aligned}
\]
\[
\left\{\tg^2l_0''+\tg l_1\tg l_2\right\}\left\{\pi-\psi-{e^2\over2}\cos l_0''(\varphi_2-\varphi_1)\right\}^2
=\left\{\tg l_1+\tg l_2\right\}^2
\]
Atqui $l_0$ non multum differt a $\pi-l_0''$; ergo summa $\cos\varphi_1+\cos\varphi_2$ minima est, quare aut $\varphi_2+\varphi_1$ aut $\varphi_2-\varphi_1$ non longe recedit a $\pi$. Quod cum in summam cadere nequeat, siquidem simplex angulus $\varphi_2$ major est quam $\pi$ (quippe longitudinem $\psi$ posuimus prope $\pi$ jacentem, itaque.
\[
\psi>\int_{l_1}^{l_2} dl\,{\cos l_0\over\cos l}\,{\sqrt{1-e^2\cos^2l}\over\sqrt{\sin^2l_0-\sin^2l}}
\qquad \text{cf. p. 35}
\]
esse) differentia \gq{$\varphi_2-\varphi_1$} prope $\pi$ sita est, ita ut loco $\varphi_2-\varphi_1$ ponere liceat $\pi$.
\endgroup
